% Luc Destombe --- licence CC BY-NC-SA 4.0
% https://creativecommons.org/licenses/by-nc-sa/4.0/deed.fr
% Fichier autonome : compiler avec pdflatex, deux fois (signets, nombre de pages).
\documentclass[11pt,a4paper]{article}
\makeatletter
\newif\iflycee@unecolonne
\lycee@unecolonnetrue
\newif\iflycee@palatino
\lycee@palatinotrue

\RequirePackage[utf8]{inputenc}
\RequirePackage[T1]{fontenc}
\RequirePackage[french]{babel}
\RequirePackage{etoolbox}

\newcommand{\ShorthandsOff}{\@ifundefined{shorthandoff}{}{\shorthandoff{;:!?}}}
\newcommand{\ShorthandsOn}{\@ifundefined{shorthandon}{}{\shorthandon{;:!?}}}
\ShorthandsOff
\AtBeginDocument{\ShorthandsOn}
\AtBeginEnvironment{tikzpicture}{\ShorthandsOff}

\RequirePackage{geometry}
\RequirePackage{fancyhdr}
\RequirePackage{lastpage}
\RequirePackage{multicol}
\RequirePackage{array}
\RequirePackage{enumitem}
\RequirePackage{xcolor}
\RequirePackage{needspace}

\RequirePackage{amsmath,amssymb}
\RequirePackage{mathpazo}
\renewcommand{\ttdefault}{lmtt}
\DeclareMathSizes{10.95}{10.95}{8}{5.5}
\begingroup\catcode`\_=\active \catcode`\^=\active
\gdef_#1{\sb{\scriptscriptstyle #1}}
\gdef^#1{\sp{\scriptscriptstyle\lycee@moinsepais #1}}
\endgroup
\newcommand\lycee@moins{\mathbin{%
  \setbox\z@\hbox{$\m@th\scriptscriptstyle\lycee@moinsorig$}%
  \dimen@=.55\wd\z@ \advance\dimen@-.5pt
  \hbox to.75\wd\z@{\hss$\m@th\scriptscriptstyle\vcenter{\hbox{%
    \tikz\draw[line width=.5pt,line cap=round](0,0)--(\the\dimen@,0);}}$\hss}}}
\begingroup\lccode`\~=`\- \lowercase{\endgroup\let~}\lycee@moins
\newcommand\lycee@moinsepais{\mathcode`\-="8000 }
\AtBeginDocument{\mathchardef\lycee@moinsorig=\mathcode`\-\relax}
\AtBeginDocument{\catcode`\_=12 \mathcode`\_="8000
  \catcode`\^=12 \mathcode`\^="8000 }
\newcommand\lycee@exposants{%
  \fontdimen13\textfont2=.48\fontdimen6\textfont2
  \fontdimen14\textfont2=.48\fontdimen6\textfont2
  \fontdimen15\textfont2=.48\fontdimen6\textfont2
  \fontdimen13\scriptfont2=.48\fontdimen6\scriptfont2
  \fontdimen14\scriptfont2=.48\fontdimen6\scriptfont2
  \fontdimen15\scriptfont2=.48\fontdimen6\scriptfont2}
\every@math@size\expandafter{\the\every@math@size\lycee@exposants}
\newlength{\retraitcalculs}
\setlength{\retraitcalculs}{2em}

\RequirePackage{tikz}
\usetikzlibrary{arrows.meta,calc,babel,shapes.misc}
\RequirePackage[most]{tcolorbox}
\tcbuselibrary{skins,breakable}

\definecolor{coulDef}{HTML}{1F4E79}
\definecolor{coulProp}{HTML}{7030A0}
\definecolor{coulRem}{HTML}{C55A11}
\definecolor{coulEx}{HTML}{2E7D32}
\definecolor{coulExo}{HTML}{B03A2E}
\definecolor{coulPart}{HTML}{1F4E79}
\definecolor{coulMeth}{HTML}{1B6E4F}
\definecolor{coulCourbe}{HTML}{0B5ED7}
\definecolor{coulCourbeB}{HTML}{C00000}
\definecolor{coulCourbeC}{HTML}{2E7D32}
\definecolor{coulGrille}{HTML}{8C8C8C}
\definecolor{coulRep}{HTML}{1B6E4F}
\definecolor{coulBar}{HTML}{2E7D32}

\newcommand{\R}{\mathbb{R}}
\newcommand{\Rs}{\mathbb{R}^{*}}
\renewcommand{\le}{\leqslant}
\renewcommand{\ge}{\geqslant}
\renewcommand{\approx}{\simeq}
\newcommand{\vect}[1]{\overrightarrow{#1}}
\newcommand{\norme}[1]{\left\lVert #1 \right\rVert}
\newcommand{\intervalle}[2]{\left[#1\,;#2\right]}
\RequirePackage{cancel}
\renewcommand{\CancelColor}{\color{coulExo}}
\newcommand{\fois}[1]{\times{\color{coulExo}#1}}
\newcommand{\barre}[1]{\cancel{#1}}

\tikzset{grille/.style={coulGrille,line width=0.4pt}}
\newcommand{\quadrillage}[4]{\draw[grille,step=1] (#1,#2) grid (#3,#4);}
\tikzset{
  croix/.style={cross out,draw,line width=0.6pt,inner sep=0pt,minimum size=4pt},
  croixdroite/.style={croix,rotate=45,line width=0.8pt},
}
\newcommand{\pt}[3]{\path (#1) node[croix]{} node[#2,font=\small,inner sep=2.5pt]{$#3$};}

\newcommand{\lycee@repere}[4]{%
  \draw[grille,step=1] (#1,#2) grid (#3,#4);
  \draw[-{Stealth[length=2mm]},thick] (#1,0)--({#3+0.4},0) node[below left=-1pt]{$x$};
  \draw[-{Stealth[length=2mm]},thick] (0,#2)--(0,{#4+0.4}) node[below left=-1pt]{$y$};
  \node[below left,font=\scriptsize,inner sep=1.5pt] at (0,0) {$0$};}
\newcommand{\lycee@gradx}[1]{\foreach \k/\l in {#1}{%
  \draw (\k,0.07)--(\k,-0.07) node[below,font=\scriptsize,inner sep=1.5pt]{$\l$};}}
\newcommand{\lycee@grady}[1]{\foreach \k/\l in {#1}{%
  \draw (0.07,\k)--(-0.07,\k) node[left,font=\scriptsize,inner sep=1.5pt]{$\l$};}}
\newcommand{\lycee@intff}[2]{\left[#1\,;#2\right]}
\newcommand{\lycee@intfo}[2]{\left[#1\,;#2\right[}
\newcommand{\lycee@intof}[2]{\left]#1\,;#2\right]}
\newcommand{\lycee@intoo}[2]{\left]#1\,;#2\right[}
\newcommand{\lycee@ens}[1]{\left\{#1\right\}}
\AtBeginDocument{%
  \providecommand{\repere}{\lycee@repere}%
  \providecommand{\gradx}{\lycee@gradx}%
  \providecommand{\grady}{\lycee@grady}%
  \providecommand{\Cf}{\mathcal{C}\sb{\scriptscriptstyle f}}%
  \providecommand{\Cg}{\mathcal{C}\sb{\scriptscriptstyle g}}%
  \providecommand{\intff}{\lycee@intff}%
  \providecommand{\intfo}{\lycee@intfo}%
  \providecommand{\intof}{\lycee@intof}%
  \providecommand{\intoo}{\lycee@intoo}%
  \providecommand{\ens}{\lycee@ens}}

\newlist{questions}{enumerate}{3}
\setlist[questions,1]{label=\arabic*\textdegree),leftmargin=*,itemsep=2pt,topsep=3pt}
\setlist[questions,2]{label=\alph*),leftmargin=*,itemsep=1pt,topsep=2pt}
\setlist[questions,3]{label=\roman*),leftmargin=*,itemsep=1pt,topsep=1pt}
\setlist[itemize]{label=\textbullet,leftmargin=*,itemsep=2pt,topsep=3pt}

\newcommand{\besoinplace}[1]{\par\penalty\@M\begingroup
  \dimen@=#1\relax\dimen@ii=\pagegoal\advance\dimen@ii-\pagetotal
  \ifdim\dimen@>\dimen@ii\ifdim\dimen@ii>\z@\vfil\fi\break\fi\endgroup}

\newcommand{\lycee@pied}{\thepage/\pageref{LastPage}}
\newcommand{\entete}[2]{%
  \pagestyle{fancy}\fancyhf{}%
  \fancyhead[L]{\footnotesize\color{coulPart}\textbf{#1}}%
  \fancyhead[R]{\footnotesize\color{coulPart}\textbf{#2}}%
  \fancyfoot[C]{\footnotesize\lycee@pied}%
  \renewcommand{\headrulewidth}{0.4pt}%
  \renewcommand{\headrule}{\hbox to\headwidth{%
    \color{coulPart!40}\leaders\hrule height \headrulewidth\hfill}}}

\RequirePackage{ccicons}
\newcommand{\lycee@licence}{{\scriptsize\color{gray}%
  \copyright\ Luc Destombe --- \ccbyncsa\ CC BY-NC-SA 4.0}}
\newif\iflycee@licence \lycee@licencetrue
\newcommand{\sanslicence}{\lycee@licencefalse}
\AtBeginDocument{\iflycee@licence\AddToHookNext{shipout/foreground}{%
  \put(0,0){\rlap{\kern\dimexpr1in+\hoffset+\oddsidemargin+\textwidth\relax
    \llap{\raisebox{-\dimexpr1in+\voffset+\topmargin+\headheight+\headsep
      +\textheight+\footskip\relax}{\lycee@licence}}}}}\fi}

\newcommand{\lycee@site}{faire-des-maths.fr}
\newcommand{\lycee@version}{{\scriptsize\color{gray}\lycee@site
  \ifdefined\lyceeversion\ --- Version~\lyceeversion\fi}}
\AtBeginDocument{\iflycee@licence\AddToHookNext{shipout/foreground}{%
  \put(0,0){\rlap{\kern\dimexpr1in+\hoffset+\oddsidemargin\relax
    \raisebox{-\dimexpr1in+\voffset+\topmargin+\headheight+\headsep
      +\textheight+\footskip\relax}{\lycee@version}}}}\fi}

\newcommand{\formulesserrees}{%
  \setlength{\abovedisplayskip}{3pt}\setlength{\belowdisplayskip}{3pt}%
  \setlength{\abovedisplayshortskip}{2pt}\setlength{\belowdisplayshortskip}{3pt}}

\setlength{\parindent}{0pt}

\RequirePackage[implicit=false,hidelinks,pdfencoding=auto,psdextra,bookmarksopen,
  bookmarksopenlevel=1,pdfcreator={},pdfproducer={}]{hyperref}
\RequirePackage{bookmark}
\hypersetup{pdfauthor={Luc Destombe},pdfkeywords={CC BY-NC-SA 4.0}}
\newcounter{lycee@signet}
\newcommand{\signet}[3][]{\stepcounter{lycee@signet}%
  \edef\lycee@dest{\if\relax\detokenize{#1}\relax
    signet.\arabic{lycee@signet}\else#1\fi}%
  \hypertarget{\lycee@dest}{}\bookmark[level=#2,dest=\lycee@dest]{#3}}
\pdfstringdefDisableCommands{%
  \def\R{R}\def\vect#1{#1}\def\Cf{Cf}\def\Cg{Cg}\def\fois#1{×#1}%
  \def\barre#1{#1}\def\intervalle#1#2{[#1;#2]}\def\intff#1#2{[#1;#2]}%
  \def\textsuperscript#1{#1}\def\dfrac#1#2{#1/#2}\def\frac#1#2{#1/#2}%
  \def\vec#1{#1⃗}}%
\RequirePackage[official]{eurosym}

\tcbset{exercice corrige/.style={
  enhanced,breakable,before upper=\raggedright,
  colback=white,colframe=coulExo,
  boxrule=0.8pt,arc=2pt,
  left=6pt,right=6pt,top=8pt,bottom=5pt,
  before skip=9pt,after skip=4pt,
  coltitle=white,fonttitle=\bfseries\small,
  attach boxed title to top left={xshift=8pt,yshift=-\tcboxedtitleheight/2},
  boxed title style={colback=coulExo,arc=2pt,boxrule=0pt}}}
\newcounter{exo}
\newtcolorbox{exercice}[1][]{exercice corrige,
  phantom={\signet[exercice-\arabic{exo}]{2}{Exercice \arabic{exo}}},
  title={Exercice \theexo\ifx\relax#1\relax\else{} --- #1\fi}}
\newenvironment{exo}[1][\relax]{\refstepcounter{exo}\begin{exercice}[#1]}{\end{exercice}}

\geometry{top=1.8cm,bottom=1cm,left=1cm,right=1cm,headheight=14pt,footskip=0.6cm}
\renewcommand{\lycee@pied}{\thepage}

\renewcommand{\theexo}{\ifnum\value{exo}<10 0\fi\arabic{exo}}
\tcbset{exercice corrige/.style={
  enhanced,breakable,before upper=\raggedright,
  colback=white,colframe=coulExo,
  boxrule=0.9pt,arc=2pt,
  left=8pt,right=8pt,top=8pt,bottom=6pt,
  before skip=8pt,after skip=6pt,
  coltitle=white,fonttitle=\bfseries,
  attach boxed title to top left={xshift=8pt,yshift=-\tcboxedtitleheight/2},
  boxed title style={colback=coulExo,arc=2pt,boxrule=0pt}}}

\newcommand{\entetefiche}[2]{%
  \begin{center}
    \begin{tcolorbox}[enhanced,width=0.92\linewidth,colback=white,colframe=coulPart,
        boxrule=1.4pt,arc=3pt,halign=center,top=6pt,bottom=6pt]
      {\large\bfseries\color{coulPart} CORRIGÉ --- FICHE D'EXERCICES}\\[3pt]
      {\normalsize\bfseries #1}\\[2pt]
      {\normalsize #2}
    \end{tcolorbox}
  \end{center}}

\newcounter{partie}
\newcommand{\partie}[1]{%
  \stepcounter{partie}%
  \besoinplace{8\baselineskip}\vspace{4pt}%
  \begin{tcolorbox}[enhanced,colback=coulPart!8,colframe=coulPart,boxrule=0pt,
      leftrule=3pt,arc=0pt,left=6pt,right=4pt,top=3pt,bottom=3pt,
      before skip=6pt,after skip=2pt]
    \bfseries\color{coulPart}\Roman{partie} --- #1\signet{1}{\Roman{partie} --- #1}
  \end{tcolorbox}}
\newcommand{\partiesansnum}[1]{%
  \besoinplace{8\baselineskip}\vspace{4pt}%
  \begin{tcolorbox}[enhanced,colback=coulPart!8,colframe=coulPart,boxrule=0pt,
      leftrule=3pt,arc=0pt,left=6pt,right=4pt,top=3pt,bottom=3pt,
      before skip=6pt,after skip=2pt]
    \bfseries\color{coulPart}#1\signet{1}{#1}
  \end{tcolorbox}}

\setlist[questions,1]{label=\arabic*\textdegree),leftmargin=*,itemsep=3pt,topsep=3pt}
\setlist[questions,2]{label=\alph*),leftmargin=*,itemsep=2pt,topsep=2pt}
\setlist[questions,3]{label=\roman*),leftmargin=*,itemsep=1pt,topsep=1pt}
\newlist{expr}{enumerate}{1}
\setlist[expr]{label={\Alph*\ $=$},leftmargin=*,itemsep=2pt,topsep=3pt}

\newcommand{\res}[1]{{\color{coulRep}#1}}
\newcommand{\rem}[1]{\par\smallskip{\small\itshape\color{black!60}$\rightarrow$ #1}}
\newenvironment{calculs}{\par\smallskip\noindent$\displaystyle\begin{aligned}[t]}%
  {\end{aligned}$\par\smallskip}
\newcommand{\justif}[1]{\quad\text{\small\itshape\color{black!60}#1}}
\newcommand{\cas}[1]{\par\smallskip\noindent\textbf{\color{coulExo}#1}\vadjust{\nobreak}}
\newcommand{\explique}[2]{\par\smallskip\noindent
  \begin{minipage}[t]{0.6\linewidth}\raggedright #1\end{minipage}\hfill
  \begin{minipage}[t]{0.37\linewidth}\raggedright\leavevmode
    {\small\itshape\color{black!60}#2}\end{minipage}\par}

\setlength{\parskip}{2pt}

\makeatother
\entete{Chapitre 10 --- Fonction exponentielle --- corrigé}{Première spécialité}
\usepackage{tkz-tab}
\let\casligne\cas
\renewcommand{\cas}[1]{\casligne{#1}\par\nopagebreak}
\setlength{\parindent}{0pt}
\newcommand{\e}{\mathrm{e}}

\begin{document}

\entetefiche{Chapitre 10 --- Fonction exponentielle}{Première spécialité mathématiques}

\partie{Définition, relation fonctionnelle, notation $\e^{x}$}

\begin{exo}[Écrire sous la forme $\e^{n}$]
\cas{a)}
\begin{calculs}
A &= \e^{3}\times\e^{6} &&\\
A &= \e^{3+6} &&\justif{$\e^{a}\times\e^{b}=\e^{a+b}$}\\
A &= \res{\e^{9}} &&
\end{calculs}
\cas{b)}
\begin{calculs}
B &= \e^{2}\times\e^{-5} &&\\
B &= \e^{2+(-5)} &&\justif{on additionne les exposants}\\
B &= \res{\e^{-3}} &&
\end{calculs}
\cas{c)}
\begin{calculs}
C &= \dfrac{\e^{9}}{\e^{4}} &&\\
C &= \e^{9-4} &&\justif{$\frac{\e^{a}}{\e^{b}}=\e^{a-b}$}\\
C &= \res{\e^{5}} &&
\end{calculs}
\cas{d)}
\begin{calculs}
D &= \dfrac{\e^{-1}}{\e^{-6}} &&\\
D &= \e^{-1-(-6)} &&\justif{on soustrait l'exposant du bas, entre parenthèses}\\
D &= \res{\e^{5}} &&\justif{$-1+6=5$}
\end{calculs}
\cas{e)}
\begin{calculs}
E &= \left(\e^{4}\right)^{3} &&\\
E &= \e^{4\times 3} &&\justif{$\left(\e^{a}\right)^{n}=\e^{na}$}\\
E &= \res{\e^{12}} &&
\end{calculs}
\cas{f)}
\begin{calculs}
F &= \dfrac{1}{\e^{-7}} &&\\
F &= \e^{-(-7)} &&\justif{$\frac{1}{\e^{a}}=\e^{-a}$}\\
F &= \res{\e^{7}} &&
\end{calculs}
\cas{g)}
\begin{calculs}
G &= \left(\e^{-2}\right)^{4}\times\e^{5} &&\\
G &= \e^{-8}\times\e^{5} &&\justif{$-2\times 4=-8$}\\
G &= \res{\e^{-3}} &&\justif{$-8+5=-3$}
\end{calculs}
\rem{Ce sont les règles des puissances de $10$ : on additionne les exposants dans un
produit, on les soustrait dans un quotient, on les multiplie pour une puissance.}
\end{exo}

\begin{exo}[Simplifier une expression]
\cas{a)}
\begin{calculs}
A &= \e^{2x}\times\e^{x+3} &&\\
A &= \e^{2x+x+3} &&\justif{$\e^{a}\times\e^{b}=\e^{a+b}$}\\
A &= \res{\e^{3x+3}} &&
\end{calculs}
\cas{b)}
\begin{calculs}
B &= \dfrac{\e^{4x-1}}{\e^{x+2}} &&\\
B &= \e^{4x-1-(x+2)} &&\justif{parenthèses autour de l'exposant du bas}\\
B &= \e^{4x-1-x-2} &&\\
B &= \res{\e^{3x-3}} &&
\end{calculs}
\cas{c)}
\begin{calculs}
C &= \left(\e^{x}\right)^{3}\times\e^{-2x} &&\\
C &= \e^{3x}\times\e^{-2x} &&\justif{$\left(\e^{a}\right)^{n}=\e^{na}$}\\
C &= \res{\e^{x}} &&\justif{$3x-2x=x$}
\end{calculs}
\cas{d)}
\begin{calculs}
D &= \dfrac{\e^{x+5}}{\e^{3x-1}} &&\\
D &= \e^{x+5-(3x-1)} &&\\
D &= \e^{x+5-3x+1} &&\justif{le « $-$ » change les deux signes}\\
D &= \res{\e^{-2x+6}} &&
\end{calculs}
\cas{e)}
\begin{calculs}
E &= \left(\e^{-x}\right)^{2}\times\e^{3x} &&\\
E &= \e^{-2x}\times\e^{3x} &&\\
E &= \res{\e^{x}} &&
\end{calculs}
\cas{f)}
\begin{calculs}
F &= \dfrac{\e^{3x}}{\e^{x}\times\e^{2x-4}} &&\\
F &= \dfrac{\e^{3x}}{\e^{3x-4}} &&\justif{d'abord le dénominateur : $x+2x-4$}\\
F &= \e^{3x-(3x-4)} &&\\
F &= \res{\e^{4}} &&\justif{$F$ ne dépend pas de $x$}
\end{calculs}
\rem{Erreur fréquente : oublier les parenthèses autour de l'exposant du dénominateur
(en b), $4x-1-x+2$ est faux).}
\end{exo}

\begin{exo}[Développer]
\cas{a)}
\begin{calculs}
A &= \left(\e^{x}+1\right)^{2} &&\\
A &= \left(\e^{x}\right)^{2}+2\times\e^{x}\times 1+1^{2} &&\justif{$(a+b)^{2}=a^{2}+2ab+b^{2}$}\\
A &= \res{\e^{2x}+2\e^{x}+1} &&\justif{$\left(\e^{x}\right)^{2}=\e^{2x}$}
\end{calculs}
\cas{b)}
\begin{calculs}
B &= \left(\e^{x}-3\right)^{2} &&\\
B &= \left(\e^{x}\right)^{2}-2\times\e^{x}\times 3+3^{2} &&\justif{$(a-b)^{2}=a^{2}-2ab+b^{2}$}\\
B &= \res{\e^{2x}-6\e^{x}+9} &&
\end{calculs}
\cas{c)}
\begin{calculs}
C &= \left(\e^{x}+2\right)\left(\e^{x}-2\right) &&\\
C &= \left(\e^{x}\right)^{2}-2^{2} &&\justif{$(a+b)(a-b)=a^{2}-b^{2}$}\\
C &= \res{\e^{2x}-4} &&
\end{calculs}
\cas{d)}
\begin{calculs}
D &= \e^{x}\left(\e^{x}+\e^{-x}\right) &&\\
D &= \e^{x}\times\e^{x}+\e^{x}\times\e^{-x} &&\justif{on distribue}\\
D &= \res{\e^{2x}+1} &&\justif{$\e^{x}\times\e^{-x}=\e^{0}=1$}
\end{calculs}
\cas{e)}
\begin{calculs}
E &= \left(\e^{x}+\e^{-x}\right)^{2} &&\\
E &= \left(\e^{x}\right)^{2}+2\times\e^{x}\times\e^{-x}+\left(\e^{-x}\right)^{2} &&\justif{$(a+b)^{2}$}\\
E &= \res{\e^{2x}+2+\e^{-2x}} &&\justif{le double produit vaut $2\times 1$}
\end{calculs}
\end{exo}

\begin{exo}[Vrai ou faux]
\cas{a)}
\explique{\res{Faux} : dans un produit, on additionne les exposants, donc
$\e^{2}\times\e^{3}=\e^{5}$.}{$\e^{a}\times\e^{b}=\e^{a+b}$}
\cas{b)}
\explique{\res{Faux} : $\e^{-x}=\dfrac{1}{\e^{x}}$ est strictement positif, alors que
$-\e^{x}$ est strictement négatif. Pour $x=0$ : $\e^{0}=1$ et $-\e^{0}=-1$.}{contre-exemple}
\cas{c)}
\explique{\res{Vrai} : une puissance d'une exponentielle multiplie l'exposant,
$\left(\e^{x}\right)^{2}=\e^{2x}$.}{$\left(\e^{a}\right)^{n}=\e^{na}$}
\cas{d)}
\explique{\res{Faux} : $\e^{x}+\e^{x}=2\e^{x}$ ; pour $x=1$, $2\e\approx 5{,}4$ et
$\e^{2}\approx 7{,}4$.}{une somme n'est pas un produit}
\cas{e)}
\explique{\res{Faux} : $\e^{0}=1$, c'est la définition ($\exp(0)=1$).}{définition 1}
\cas{f)}
\explique{\res{Vrai} : c'est la relation $\exp(x)\times\exp(-x)=1$ du cours.}{propriété 1}
\end{exo}

\begin{exo}[Une expression constante]
\cas{1.}
\begin{calculs}
f(0) &= \left(\e^{0}+\e^{0}\right)^{2}-\left(\e^{0}-\e^{0}\right)^{2} &&\\
f(0) &= 2^{2}-0^{2} &&\justif{$\e^{0}=1$}\\
f(0) &= \res{4} &&
\end{calculs}
\explique{À la calculatrice, $f(1)=\res{4}$ aussi.}{on conjecture que $f(x)$ vaut toujours $4$}
\cas{2.}
\begin{calculs}
\left(\e^{x}+\e^{-x}\right)^{2} &= \res{\e^{2x}+2+\e^{-2x}} &&\justif{exercice 3 e)}\\
\left(\e^{x}-\e^{-x}\right)^{2} &= \left(\e^{x}\right)^{2}-2\times\e^{x}\times\e^{-x}+\left(\e^{-x}\right)^{2} &&\justif{$(a-b)^{2}$}\\
\left(\e^{x}-\e^{-x}\right)^{2} &= \res{\e^{2x}-2+\e^{-2x}} &&\justif{$\e^{x}\times\e^{-x}=1$}
\end{calculs}
\cas{3.}
\begin{calculs}
f(x) &= \left(\e^{2x}+2+\e^{-2x}\right)-\left(\e^{2x}-2+\e^{-2x}\right) &&\\
f(x) &= \e^{2x}+2+\e^{-2x}-\e^{2x}+2-\e^{-2x} &&\justif{le « $-$ » change tous les signes}\\
f(x) &= \res{4} &&
\end{calculs}
\explique{Pour tout nombre réel $x$, $f(x)=4$ : $f(x)$ ne dépend pas de $x$.}{la conjecture
de la question 1 est démontrée}
\end{exo}

\partie{Équations et inéquations avec des exponentielles}

\begin{exo}[Résoudre une équation]
\cas{a)}
\begin{calculs}
\e^{3x+1} &= \e^{x+7} &&\\
3x+1 &= x+7 &&\justif{$\e^{a}=\e^{b}\iff a=b$}\\
2x &= 6 &&\\
x &= 3 &&
\end{calculs}
\explique{$S=\res{\ens{3}}$}{}
\cas{b)}
\begin{calculs}
\e^{5x-10} &= \e^{0} &&\justif{$1=\e^{0}$}\\
5x-10 &= 0 &&\\
x &= 2 &&
\end{calculs}
\explique{$S=\res{\ens{2}}$}{}
\cas{c)}
\begin{calculs}
\e^{2x+3} &= \e^{1-x} &&\\
2x+3 &= 1-x &&\\
3x &= -2 &&\\
x &= -\dfrac{2}{3} &&
\end{calculs}
\explique{$S=\res{\ens{-\frac{2}{3}}}$}{}
\cas{d)}
\begin{calculs}
\e^{x} &= \e^{1} &&\justif{$\e=\e^{1}$}\\
x &= 1 &&
\end{calculs}
\explique{$S=\res{\ens{1}}$}{}
\cas{e)}
\explique{$\e^{4-x}>0$ pour tout nombre réel $x$ : il ne vaut jamais $-2$, donc
$S=\res{\varnothing}$.}{une exponentielle est strictement positive}
\end{exo}

\begin{exo}[Équations et second degré]
\cas{a)}
\begin{calculs}
\e^{x^{2}} &= \e^{2x+8} &&\\
x^{2} &= 2x+8 &&\justif{on égale les exposants}\\
x^{2}-2x-8 &= 0 &&\\
\Delta &= (-2)^{2}-4\times 1\times(-8)=36 &&\\
x_{1} &= \dfrac{2-6}{2}=-2 &&\justif{$\frac{-b-\sqrt{\Delta}}{2a}$}\\
x_{2} &= \dfrac{2+6}{2}=4 &&
\end{calculs}
\explique{$S=\res{\ens{-2\,;4}}$}{}
\cas{b)}
\begin{calculs}
\e^{x^{2}-5x} &= \e^{0} &&\justif{$1=\e^{0}$}\\
x^{2}-5x &= 0 &&\\
x(x-5) &= 0 &&\justif{facteur commun $x$ : inutile de calculer $\Delta$}
\end{calculs}
\explique{Un produit est nul si l'un de ses facteurs est nul : $S=\res{\ens{0\,;5}}$.}{$x=0$
ou $x-5=0$}
\cas{c)}
\begin{calculs}
x^{2}+1 &= x &&\justif{on égale les exposants}\\
x^{2}-x+1 &= 0 &&\\
\Delta &= (-1)^{2}-4\times 1\times 1=-3 &&
\end{calculs}
\explique{$\Delta<0$ : pas de solution, $S=\res{\varnothing}$.}{}
\cas{d)}
\begin{calculs}
\e^{2x}\times\e^{x^{2}} &= \e^{3} &&\\
\e^{x^{2}+2x} &= \e^{3} &&\justif{$\e^{a}\times\e^{b}=\e^{a+b}$}\\
x^{2}+2x-3 &= 0 &&\\
\Delta &= 2^{2}-4\times 1\times(-3)=16 &&\\
x_{1} &= \dfrac{-2-4}{2}=-3 &&\\
x_{2} &= \dfrac{-2+4}{2}=1 &&
\end{calculs}
\explique{$S=\res{\ens{-3\,;1}}$}{}
\end{exo}

\begin{exo}[Résoudre une inéquation]
\cas{a)}
\begin{calculs}
\e^{2x+5} &\le \e^{x+8} &&\\
2x+5 &\le x+8 &&\justif{$\exp$ est croissante : même sens}\\
x &\le 3 &&
\end{calculs}
\explique{$S=\res{\intof{-\infty}{3}}$}{}
\cas{b)}
\begin{calculs}
\e^{3-x} &> \e^{0} &&\justif{$1=\e^{0}$}\\
3-x &> 0 &&\\
-x &> -3 &&\\
x &< 3 &&\justif{on multiplie par $-1<0$ : le sens change}
\end{calculs}
\explique{$S=\res{\intoo{-\infty}{3}}$}{}
\cas{c)}
\begin{calculs}
\e^{4x} &\ge \e^{x-6} &&\\
4x &\ge x-6 &&\\
3x &\ge -6 &&\\
x &\ge -2 &&\justif{on divise par $3>0$ : même sens}
\end{calculs}
\explique{$S=\res{\intfo{-2}{+\infty}}$}{}
\cas{d)}
\begin{calculs}
x^{2} &< x+2 &&\justif{même sens}\\
x^{2}-x-2 &< 0 &&\\
\Delta &= (-1)^{2}-4\times 1\times(-2)=9 &&\\
x_{1} &= \dfrac{1-3}{2}=-1 &&\\
x_{2} &= \dfrac{1+3}{2}=2 &&
\end{calculs}
\explique{Le trinôme est négatif entre ses racines : $S=\res{\intoo{-1}{2}}$.}{signe de $a=1$
à l'extérieur des racines}
\cas{e)}
\begin{calculs}
x^{2} &\ge 4 &&\justif{même sens}\\
x^{2}-4 &\ge 0 &&\\
(x-2)(x+2) &\ge 0 &&\justif{racines $-2$ et $2$}
\end{calculs}
\explique{Le trinôme est positif à l'extérieur de ses racines :
$S=\res{\intof{-\infty}{-2}\cup\intfo{2}{+\infty}}$.}{$a=1>0$}
\rem{Erreur fréquente en e) : écrire $x\ge 2$ seulement ; $(-3)^{2}=9\ge 4$ aussi.}
\end{exo}

\begin{exo}[Produit par une exponentielle]
\cas{a)}
\explique{$\e^{x}\ne 0$ : le produit est nul quand $x-3=0$, donc $S=\res{\ens{3}}$.}{produit
nul}
\cas{b)}
\explique{$\e^{x}>0$ : le produit a le signe de $2x+1$.}{on ne garde que le facteur qui peut
changer de signe}
\begin{calculs}
2x+1 &> 0 &&\\
x &> -\dfrac{1}{2} &&
\end{calculs}
\explique{$S=\res{\intoo{-\frac{1}{2}}{+\infty}}$}{}
\cas{c)}
\explique{$\e^{-x}>0$ : le produit a le signe de $x^{2}-4=(x-2)(x+2)$, négatif entre ses
racines.}{trinôme, $a=1>0$}
\explique{$S=\res{\intff{-2}{2}}$}{les bornes conviennent : $\le 0$}
\end{exo}

\begin{exo}[Changement de variable]
\cas{1.}
\begin{calculs}
\e^{2x} &= \left(\e^{x}\right)^{2} &&\justif{$\e^{na}=\left(\e^{a}\right)^{n}$}\\
\e^{2x} &= \res{X^{2}} &&\justif{$X=\e^{x}$}
\end{calculs}
\cas{2.}
\begin{calculs}
X^{2}+2X-3 &= 0 &&\\
\Delta &= 2^{2}-4\times 1\times(-3)=16 &&\\
X_{1} &= \dfrac{-2-4}{2}=-3 &&\\
X_{2} &= \dfrac{-2+4}{2}=1 &&
\end{calculs}
\explique{Les solutions sont \res{$X=-3$ et $X=1$}.}{}
\cas{3.}
\explique{$\e^{x}=-3$ est impossible, car une exponentielle est strictement positive.}{on
revient à $x$ : $X=\e^{x}$}
\begin{calculs}
\e^{x} &= 1 &&\\
\e^{x} &= \e^{0} &&\\
x &= 0 &&
\end{calculs}
\explique{$S=\res{\ens{0}}$}{}
\end{exo}

\partie{Étude de la fonction exponentielle}

\begin{exo}[Dériver une fonction qui contient $\e^{x}$]
\cas{a)}
\begin{calculs}
f(x) &= 4\e^{x}+x^{2} &&\\
f'(x) &= \res{4\e^{x}+2x} &&\justif{$\left(\e^{x}\right)'=\e^{x}$ ; $\left(x^{2}\right)'=2x$}
\end{calculs}
\cas{b)}
\begin{calculs}
u(x) &= x+3 &&\justif{$u'(x)=1$}\\
v(x) &= \e^{x} &&\justif{$v'(x)=\e^{x}$}\\
g'(x) &= 1\times\e^{x}+(x+3)\,\e^{x} &&\justif{$u'v+uv'$}\\
g'(x) &= \res{(x+4)\,\e^{x}} &&\justif{on factorise par $\e^{x}$ : $1+x+3$}
\end{calculs}
\cas{c)}
\begin{calculs}
u(x) &= 3-x &&\justif{$u'(x)=-1$}\\
h'(x) &= -1\times\e^{x}+(3-x)\,\e^{x} &&\justif{$u'v+uv'$}\\
h'(x) &= \res{(2-x)\,\e^{x}} &&\justif{$-1+3-x$}
\end{calculs}
\cas{d)}
\begin{calculs}
u(x) &= \e^{x} &&\justif{$u'(x)=\e^{x}$}\\
v(x) &= x &&\justif{$v'(x)=1$}\\
k'(x) &= \dfrac{\e^{x}\times x-\e^{x}\times 1}{x^{2}} &&\justif{$\frac{u'v-uv'}{v^{2}}$}\\
k'(x) &= \res{\dfrac{(x-1)\,\e^{x}}{x^{2}}} &&\justif{on factorise le numérateur par $\e^{x}$}
\end{calculs}
\cas{e)}
\begin{calculs}
u(x) &= x^{2} &&\justif{$u'(x)=2x$}\\
m'(x) &= 2x\,\e^{x}+x^{2}\e^{x} &&\justif{$u'v+uv'$, avec $v(x)=\e^{x}$}\\
m'(x) &= \res{x(x+2)\,\e^{x}} &&\justif{$\left(x^{2}+2x\right)\e^{x}$, puis facteur $x$}
\end{calculs}
\end{exo}

\begin{exo}[Tangentes à la courbe de l'exponentielle]
\cas{1.}
\explique{Au point d'abscisse $a$, la tangente a pour équation $y=\e^{a}(x-a)+\e^{a}$.}{$y=f'(a)(x-a)+f(a)$,
avec $f'=f=\exp$}
\begin{calculs}
y &= \e^{-1}(x+1)+\e^{-1} &&\justif{$a=-1$}\\
y &= \res{\e^{-1}x+2\e^{-1}} &&\\
y &= \e^{0}(x-0)+\e^{0} &&\justif{$a=0$}\\
y &= \res{x+1} &&\\
y &= \e^{2}(x-2)+\e^{2} &&\justif{$a=2$}\\
y &= \res{\e^{2}x-\e^{2}} &&\justif{$-2\e^{2}+\e^{2}=-\e^{2}$}
\end{calculs}
\cas{2.}
\explique{On résout $y=0$ dans chaque équation.}{point de l'axe des abscisses : $y=0$}
\begin{calculs}
\e^{-1}x+2\e^{-1} &= 0 &&\justif{$a=-1$}\\
\e^{-1}(x+2) &= 0 &&\\
x &= \res{-2} &&\justif{$\e^{-1}\ne 0$}\\
x+1 &= 0 &&\justif{$a=0$}\\
x &= \res{-1} &&\\
\e^{2}(x-1) &= 0 &&\justif{$a=2$}\\
x &= \res{1} &&
\end{calculs}
\cas{3.}
\explique{\res{La tangente au point d'abscisse $a$ coupe l'axe des abscisses au point
d'abscisse $a-1$.}}{$-1-1=-2$ ; $0-1=-1$ ; $2-1=1$}
\rem{En effet, $\e^{a}(x-a)+\e^{a}=\e^{a}(x-a+1)$ s'annule pour $x=a-1$. Cela donne une
façon rapide de tracer une tangente à la courbe de l'exponentielle.}
\end{exo}

\begin{exo}[Courbe et tangente]
\cas{1.}
\begin{calculs}
\varphi(x) &= \e^{x}-x-1 &&\\
\varphi'(x) &= \res{\e^{x}-1} &&\justif{$\left(\e^{x}\right)'=\e^{x}$ ; $(x)'=1$ ; $(1)'=0$}
\end{calculs}
\cas{2.}
\begin{calculs}
\e^{x}-1 &\ge 0 &&\\
\e^{x} &\ge \e^{0} &&\justif{$1=\e^{0}$}\\
x &\ge 0 &&\justif{$\exp$ est croissante}
\end{calculs}
\explique{$S=\res{\intfo{0}{+\infty}}$ : $\varphi'(x)$ est négatif pour $x<0$, positif pour
$x>0$.}{}
\cas{3.}
\begin{calculs}
\varphi(0) &= \e^{0}-0-1 &&\\
\varphi(0) &= 0 &&
\end{calculs}
\begin{center}
\begin{tikzpicture}[scale=0.9]
\tkzTabInit[lgt=2,espcl=2.4]{$x$/0.8,$\varphi'(x)$/0.8,$\varphi(x)$/1.6}{$-\infty$,$0$,$+\infty$}
\tkzTabLine{,-,z,+,}
\tkzTabVar{+/{},-/$0$,+/{}}
\end{tikzpicture}
\end{center}
\explique{Le minimum de $\varphi$ est \res{$0$, atteint en $0$}.}{$\varphi'$ s'annule en changeant
de signe}
\cas{4.}
\explique{Le minimum est $0$ : pour tout nombre réel $x$, $\varphi(x)\ge 0$, soit
$\e^{x}-x-1\ge 0$, donc \res{$\e^{x}\ge x+1$}.}{on ajoute $x+1$ aux deux membres}
\cas{5.}
\explique{$y=x+1$ est la tangente à la courbe de l'exponentielle au point d'abscisse $0$
(exercice 12) : \res{la courbe est toujours au-dessus de cette tangente}.}{elles se touchent
au point $(0\,;1)$}
\end{exo}

\begin{exo}[Comparer sans calculatrice]
\cas{1.}
\explique{La fonction exponentielle est strictement croissante : les exponentielles sont
rangées dans le même ordre que leurs exposants.}{propriété 4}
\explique{$0{,}3<0{,}5$, donc \res{$\e^{0{,}3}<\e^{0{,}5}$} ; $-3<-2$, donc
\res{$\e^{-3}<\e^{-2}$} ; $-0{,}1<0$, donc \res{$\e^{-0{,}1}<1$}.}{$1=\e^{0}$}
\cas{2.}
\explique{$-1\le x\le 2$ et $\exp$ est croissante : \res{$\e^{-1}\le\e^{x}\le\e^{2}$}.}{environ
$0{,}37\le\e^{x}\le 7{,}39$}
\cas{3.}
\explique{$\e^{x}-1$ est \res{négatif pour $x<0$, nul pour $x=0$, positif pour $x>0$}.}{on
compare $\e^{x}$ à $\e^{0}=1$ (exercice 13)}
\end{exo}

\partie{Dérivée de $\e^{u}$}

\begin{exo}[Dériver $\e^{u}$]
\cas{a)}
\begin{calculs}
u(x) &= 3x &&\justif{$u'(x)=3$}\\
f'(x) &= \res{3\e^{3x}} &&\justif{$\left(\e^{u}\right)'=u'\e^{u}$}
\end{calculs}
\cas{b)}
\begin{calculs}
u(x) &= -x &&\justif{$u'(x)=-1$}\\
g'(x) &= \res{-\e^{-x}} &&
\end{calculs}
\cas{c)}
\begin{calculs}
u(x) &= 2x-5 &&\justif{$u'(x)=2$}\\
h'(x) &= \res{2\e^{2x-5}} &&
\end{calculs}
\cas{d)}
\begin{calculs}
k'(t) &= 50\times(-0{,}2)\,\e^{-0{,}2t} &&\justif{$\left(\e^{at}\right)'=a\,\e^{at}$}\\
k'(t) &= \res{-10\,\e^{-0{,}2t}} &&
\end{calculs}
\cas{e)}
\begin{calculs}
u(x) &= x^{2}+1 &&\justif{$u'(x)=2x$}\\
m'(x) &= \res{2x\,\e^{x^{2}+1}} &&
\end{calculs}
\cas{f)}
\begin{calculs}
u(x) &= 1-x^{2} &&\justif{$u'(x)=-2x$}\\
p'(x) &= \res{-2x\,\e^{1-x^{2}}} &&
\end{calculs}
\rem{L'exposant reste tel quel : la dérivée de $\e^{x^{2}+1}$ n'est pas $\e^{2x}$.}
\end{exo}

\begin{exo}[Produit et exponentielle]
\cas{a)}
\begin{calculs}
u(x) &= x &&\justif{$u'(x)=1$}\\
v(x) &= \e^{3x} &&\justif{$v'(x)=3\e^{3x}$}\\
f'(x) &= 1\times\e^{3x}+x\times 3\e^{3x} &&\justif{$u'v+uv'$}\\
f'(x) &= \res{(1+3x)\,\e^{3x}} &&\justif{on factorise par $\e^{3x}$}
\end{calculs}
\cas{b)}
\begin{calculs}
u(x) &= 2x+1 &&\justif{$u'(x)=2$}\\
v(x) &= \e^{-x} &&\justif{$v'(x)=-\e^{-x}$}\\
g'(x) &= 2\e^{-x}+(2x+1)\times\left(-\e^{-x}\right) &&\\
g'(x) &= (2-2x-1)\,\e^{-x} &&\justif{on factorise par $\e^{-x}$}\\
g'(x) &= \res{(1-2x)\,\e^{-x}} &&
\end{calculs}
\cas{c)}
\begin{calculs}
u(x) &= x+1 &&\justif{$u'(x)=1$}\\
v(x) &= \e^{2x} &&\justif{$v'(x)=2\e^{2x}$}\\
h'(x) &= \e^{2x}+(x+1)\times 2\e^{2x} &&\\
h'(x) &= (1+2x+2)\,\e^{2x} &&\\
h'(x) &= \res{(2x+3)\,\e^{2x}} &&
\end{calculs}
\cas{d)}
\begin{calculs}
u(x) &= x^{2} &&\justif{$u'(x)=2x$}\\
v(x) &= \e^{-x} &&\justif{$v'(x)=-\e^{-x}$}\\
k'(x) &= 2x\,\e^{-x}-x^{2}\e^{-x} &&\\
k'(x) &= \res{x(2-x)\,\e^{-x}} &&\justif{facteurs communs $x$ et $\e^{-x}$}
\end{calculs}
\end{exo}

\begin{exo}[Une dérivée proportionnelle à la fonction]
\cas{1.}
\begin{calculs}
f'(x) &= 3\times(-2)\,\e^{-2x} &&\justif{$\left(\e^{ax}\right)'=a\,\e^{ax}$}\\
f'(x) &= -6\,\e^{-2x} &&\\
f'(x)+2f(x) &= -6\,\e^{-2x}+2\times 3\,\e^{-2x} &&\\
f'(x)+2f(x) &= 0 &&
\end{calculs}
\cas{2.}
\begin{calculs}
g'(x) &= 4\e^{4x+1} &&\justif{$u(x)=4x+1$, $u'(x)=4$}\\
g'(x) &= 4g(x) &&\justif{$\e^{4x+1}=g(x)$}
\end{calculs}
\cas{3.}
\explique{Comme au 2., $\e^{5x}$ a pour dérivée $5\e^{5x}$ ; on la multiplie par $2$ pour avoir
$h(0)=2$ : \res{$h(x)=2\e^{5x}$}.}{on s'inspire des questions 1 et 2}
\begin{calculs}
h'(x) &= 10\,\e^{5x} &&\\
h'(x) &= 5\times 2\e^{5x}=5h(x) &&\justif{vérification}\\
h(0) &= 2\e^{0}=2 &&
\end{calculs}
\rem{Ce n'est pas la seule idée possible au départ, mais c'est la seule fonction qui convient
(comme pour la définition de $\exp$, admis).}
\end{exo}

\begin{exo}[Une courbe en cloche]
\cas{1.}
\begin{calculs}
u(x) &= -0{,}25x^{2} &&\justif{$u'(x)=-0{,}5x$}\\
f'(x) &= \res{-0{,}5x\,\e^{-0{,}25x^{2}}} &&\justif{$\left(\e^{u}\right)'=u'\e^{u}$}
\end{calculs}
\cas{2.}
\explique{$\e^{-0{,}25x^{2}}>0$ : $f'(x)$ a le signe de $-0{,}5x$, positif pour $x<0$ et négatif
pour $x>0$.}{$-0{,}5<0$}
\begin{calculs}
f(-3) &= \e^{-2{,}25} &&\justif{environ $0{,}11$}\\
f(0) &= \e^{0}=1 &&\\
f(3) &= \e^{-2{,}25} &&
\end{calculs}
\begin{center}
\begin{tikzpicture}[scale=0.9]
\tkzTabInit[lgt=2,espcl=2.4]{$x$/0.8,$f'(x)$/0.8,$f(x)$/1.6}{$-3$,$0$,$3$}
\tkzTabLine{,+,z,-,}
\tkzTabVar{-/$\e^{-2{,}25}$,+/$1$,-/$\e^{-2{,}25}$}
\end{tikzpicture}
\end{center}
\explique{$f$ est \res{croissante sur $\intff{-3}{0}$, décroissante sur $\intff{0}{3}$}.}{}
\cas{3.}
\explique{Le maximum de $f$ est \res{$1$, atteint en $x=0$}.}{$f'$ s'annule en changeant de
signe}
\cas{4.}
\begin{center}
\renewcommand{\arraystretch}{1.3}
\begin{tabular}{|c|c|c|c|c|c|c|c|}
\hline
$x$ & $-3$ & $-2$ & $-1$ & $0$ & $1$ & $2$ & $3$\\ \hline
$f(x)$ & $0{,}11$ & $0{,}37$ & $0{,}78$ & $1$ & $0{,}78$ & $0{,}37$ & $0{,}11$\\ \hline
\end{tabular}
\end{center}
\begin{center}
\begin{tikzpicture}[x=2cm,y=2cm]
  \repere{-3.4}{-0.4}{3.4}{1.4}
  \gradx{-3,-2,-1,1,2,3} \grady{1}
  \draw[coulCourbe,line width=1.1pt,domain=-3:3,samples=100]
    plot (\x,{exp(-0.25*\x*\x)});
  \path[coulCourbe] (-3,0.105) node[croix]{} (-2,0.368) node[croix]{} (-1,0.779) node[croix]{}
    (0,1) node[croix]{} (1,0.779) node[croix]{} (2,0.368) node[croix]{} (3,0.105) node[croix]{};
  \node[coulCourbe,font=\footnotesize] at (1.7,0.75) {$\Cf$};
\end{tikzpicture}
\end{center}
\explique{On place les points du tableau, puis on les relie ; la tangente est horizontale au
sommet $(0\,;1)$.}{la courbe est symétrique par rapport à l'axe des ordonnées}
\end{exo}

\partie{Étude de fonctions avec des exponentielles}

\begin{exo}[Étudier une fonction]
\cas{1.}
\begin{calculs}
u(x) &= 2x-3 &&\justif{$u'(x)=2$}\\
v(x) &= \e^{x} &&\justif{$v'(x)=\e^{x}$}\\
f'(x) &= 2\e^{x}+(2x-3)\,\e^{x} &&\justif{$u'v+uv'$}\\
f'(x) &= \res{(2x-1)\,\e^{x}} &&\justif{$2+2x-3$}
\end{calculs}
\cas{2.}
\explique{$\e^{x}>0$ : $f'(x)$ a le signe de $2x-1$, qui s'annule en $\frac{1}{2}$ :
\res{négatif sur $\intff{-2}{\frac{1}{2}}$, positif sur $\intff{\frac{1}{2}}{2}$}.}{fonction
affine croissante}
\cas{3.}
\begin{calculs}
f(-2) &= -7\e^{-2} &&\justif{environ $-0{,}95$}\\
f\left(\tfrac{1}{2}\right) &= (1-3)\,\e^{\frac{1}{2}}=-2\e^{\frac{1}{2}} &&\justif{environ $-3{,}30$}\\
f(2) &= \e^{2} &&\justif{environ $7{,}39$}
\end{calculs}
\begin{center}
\begin{tikzpicture}[scale=0.9]
\tkzTabInit[lgt=2,espcl=2.6]{$x$/0.8,$f'(x)$/0.8,$f(x)$/1.6}{$-2$,$\frac{1}{2}$,$2$}
\tkzTabLine{,-,z,+,}
\tkzTabVar{+/$-7\e^{-2}$,-/$-2\e^{\frac{1}{2}}$,+/$\e^{2}$}
\end{tikzpicture}
\end{center}
\explique{Le minimum de $f$ est \res{$-2\e^{\frac{1}{2}}$ (environ $-3{,}30$), atteint en
$\frac{1}{2}$}.}{}
\cas{4.}
\begin{calculs}
f(0) &= -3 &&\\
f'(0) &= (0-1)\,\e^{0}=-1 &&\\
y &= -1\times(x-0)-3 &&\justif{$y=f'(0)(x-0)+f(0)$}\\
y &= \res{-x-3} &&
\end{calculs}
\cas{5.}
\begin{center}
\renewcommand{\arraystretch}{1.3}
\begin{tabular}{|c|c|c|c|c|c|}
\hline
$x$ & $-2$ & $-1$ & $0$ & $1$ & $2$\\ \hline
$f(x)$ & $-0{,}9$ & $-1{,}8$ & $-3$ & $-2{,}7$ & $7{,}4$\\ \hline
\end{tabular}
\end{center}
\begin{center}
\begin{tikzpicture}[x=0.6cm,y=0.6cm]
  \repere{-3}{-4}{3}{8}
  \gradx{-2,-1,1,2} \grady{-3,-2,-1,1,2,3,4,5,6,7}
  \draw[coulCourbe,line width=1.1pt,domain=-2:2,samples=100]
    plot (\x,{(2*\x-3)*exp(\x)});
  \draw[coulExo,line width=0.9pt,domain=-2.2:1] plot (\x,{-\x-3});
  \path[coulCourbe] (0.5,-3.297) node[croix]{} (0,-3) node[croix]{};
  \node[coulExo,font=\footnotesize] at (-2.4,-0.4) {$T$};
  \node[coulCourbe,font=\footnotesize] at (2.4,6.5) {$\Cf$};
\end{tikzpicture}
\end{center}
\explique{On place le minimum $\left(\frac{1}{2}\,;-3{,}3\right)$, où la tangente est
horizontale, puis $T$ au point $(0\,;-3)$.}{$T$ passe aussi par $(-3\,;0)$}
\end{exo}

\begin{exo}[Avec $\e^{-x}$]
\cas{1.}
\begin{calculs}
u(x) &= x+1 &&\justif{$u'(x)=1$}\\
v(x) &= \e^{-x} &&\justif{$v'(x)=-\e^{-x}$}\\
f'(x) &= 1\times\e^{-x}+(x+1)\times\left(-\e^{-x}\right) &&\justif{$u'v+uv'$}\\
f'(x) &= (1-x-1)\,\e^{-x} &&\justif{on factorise par $\e^{-x}$}\\
f'(x) &= -x\,\e^{-x} &&
\end{calculs}
\cas{2.}
\explique{$\e^{-x}>0$ : $f'(x)$ a le signe de $-x$, positif sur $\intff{-1}{0}$, négatif sur
$\intff{0}{4}$.}{}
\begin{calculs}
f(-1) &= 0\times\e^{1}=0 &&\\
f(0) &= 1\times\e^{0}=1 &&\\
f(4) &= 5\e^{-4} &&\justif{environ $0{,}09$}
\end{calculs}
\begin{center}
\begin{tikzpicture}[scale=0.9]
\tkzTabInit[lgt=2,espcl=2.4]{$x$/0.8,$f'(x)$/0.8,$f(x)$/1.6}{$-1$,$0$,$4$}
\tkzTabLine{,+,z,-,}
\tkzTabVar{-/$0$,+/$1$,-/$5\e^{-4}$}
\end{tikzpicture}
\end{center}
\explique{$f$ est \res{croissante sur $\intff{-1}{0}$, décroissante sur $\intff{0}{4}$}.}{}
\cas{3.}
\explique{Le maximum de $f$ sur $\intff{-1}{4}$ est $f(0)=1$ : \res{l'affirmation est vraie}.}{on
lit le maximum dans le tableau}
\cas{4.}
\begin{calculs}
f(1) &= 2\e^{-1} &&\\
f'(1) &= -\e^{-1} &&\\
y &= -\e^{-1}(x-1)+2\e^{-1} &&\justif{$y=f'(1)(x-1)+f(1)$}\\
y &= \res{-\e^{-1}x+3\e^{-1}} &&\justif{$\e^{-1}+2\e^{-1}=3\e^{-1}$}
\end{calculs}
\end{exo}

\begin{exo}[Un quotient]
\cas{1.}
\begin{calculs}
u(x) &= \e^{x} &&\justif{$u'(x)=\e^{x}$}\\
v(x) &= x^{2}+1 &&\justif{$v'(x)=2x$ ; $v(x)\ge 1$ ne s'annule pas}\\
f'(x) &= \dfrac{\e^{x}\left(x^{2}+1\right)-\e^{x}\times 2x}{\left(x^{2}+1\right)^{2}} &&\justif{$\frac{u'v-uv'}{v^{2}}$}\\
f'(x) &= \dfrac{\left(x^{2}-2x+1\right)\e^{x}}{\left(x^{2}+1\right)^{2}} &&\justif{on factorise par $\e^{x}$}\\
f'(x) &= \dfrac{(x-1)^{2}\,\e^{x}}{\left(x^{2}+1\right)^{2}} &&\justif{$(a-b)^{2}=a^{2}-2ab+b^{2}$}
\end{calculs}
\cas{2.}
\explique{Un carré est positif ou nul, $\e^{x}>0$ et le dénominateur est un carré non nul :
$f'(x)\ge 0$, et $f'(x)=0$ seulement pour $x=1$.}{}
\explique{\res{$f$ est strictement croissante sur $\R$}.}{une dérivée qui ne s'annule qu'en un
point ne change rien}
\cas{3.}
\begin{calculs}
f(1) &= \dfrac{\e}{2} &&\\
f'(1) &= 0 &&\\
y &= 0\times(x-1)+\dfrac{\e}{2} &&\\
y &= \res{\dfrac{\e}{2}} &&\justif{tangente horizontale}
\end{calculs}
\cas{4.}
\explique{\res{C'est faux} : $f'(1)=0$, mais $f'$ ne change pas de signe en $1$ ; $f$ est
croissante avant et après $1$, donc n'a pas d'extremum en $1$.}{un extremum demande un
changement de signe de $f'$}
\rem{En $1$, la tangente est horizontale mais la courbe la traverse.}
\end{exo}

\begin{exo}[Tangentes horizontales]
\cas{1.}
\explique{$\e^{x}\ne 0$ : $f(x)=0$ quand $x^{2}-3=0$, soit $x^{2}=3$.}{produit nul}
\explique{$S=\res{\ens{-\sqrt{3}\,;\sqrt{3}}}$}{deux solutions opposées}
\cas{2.}
\begin{calculs}
f'(x) &= 2x\,\e^{x}+\left(x^{2}-3\right)\e^{x} &&\justif{$u'v+uv'$, $u(x)=x^{2}-3$, $v(x)=\e^{x}$}\\
f'(x) &= \left(x^{2}+2x-3\right)\e^{x} &&\justif{on factorise par $\e^{x}$}\\
f'(x) &= (x+3)(x-1)\,\e^{x} &&\justif{on développe $(x+3)(x-1)$ pour vérifier}
\end{calculs}
\cas{3.}
\explique{Tangente horizontale : $f'(x)=0$, soit $x=-3$ ou $x=1$.}{$\e^{x}\ne 0$}
\begin{calculs}
f(-3) &= (9-3)\,\e^{-3}=6\e^{-3} &&\justif{environ $0{,}30$}\\
f(1) &= (1-3)\,\e^{1}=-2\e &&\justif{environ $-5{,}44$}
\end{calculs}
\explique{Les points sont \res{$\left(-3\,;6\e^{-3}\right)$ et $(1\,;-2\e)$}.}{}
\cas{4.}
\explique{$\e^{x}>0$ : $f'(x)$ a le signe de $(x+3)(x-1)$, positif à l'extérieur de $-3$ et
$1$, négatif entre.}{trinôme, $a=1>0$}
\begin{center}
\begin{tikzpicture}[scale=0.9]
\tkzTabInit[lgt=2,espcl=2.4]{$x$/0.8,$f'(x)$/0.8,$f(x)$/1.6}{$-\infty$,$-3$,$1$,$+\infty$}
\tkzTabLine{,+,z,-,z,+,}
\tkzTabVar{-/{},+/$6\e^{-3}$,-/$-2\e$,+/{}}
\end{tikzpicture}
\end{center}
\explique{$f$ a un \res{maximum local $6\e^{-3}$ en $-3$} et un \res{minimum local $-2\e$ en
$1$}.}{$f'$ s'annule en changeant de signe}
\end{exo}

\partie{Croissance et décroissance exponentielles}

\begin{exo}[Tracer $\e^{kt}$ et $\e^{-kt}$]
\cas{1.}
\explique{$f(t)=\e^{kt}$ avec $k=0{,}3>0$ : $f$ est \res{strictement croissante} ; $g(t)=\e^{-kt}$ :
$g$ est \res{strictement décroissante}.}{$f'(t)=0{,}3\,\e^{0{,}3t}>0$ ;
$g'(t)=-0{,}3\,\e^{-0{,}3t}<0$}
\cas{2.}
\begin{center}
\renewcommand{\arraystretch}{1.3}
\begin{tabular}{|c|c|c|c|c|c|c|c|c|c|}
\hline
$t$ & $-4$ & $-3$ & $-2$ & $-1$ & $0$ & $1$ & $2$ & $3$ & $4$\\ \hline
$f(t)$ & $0{,}3$ & $0{,}4$ & $0{,}5$ & $0{,}7$ & $1$ & $1{,}3$ & $1{,}8$ & $2{,}5$ & $3{,}3$\\ \hline
$g(t)$ & $3{,}3$ & $2{,}5$ & $1{,}8$ & $1{,}3$ & $1$ & $0{,}7$ & $0{,}5$ & $0{,}4$ & $0{,}3$\\ \hline
\end{tabular}
\end{center}
\explique{La ligne de $g$ est celle de $f$ lue de droite à gauche.}{$g(t)=f(-t)$}
\par\noindent\begin{minipage}{\linewidth}
\cas{3.}
\begin{center}
\begin{tikzpicture}[x=1cm,y=1cm]
  \repere{-5}{-1}{5}{4}
  \gradx{-4,-3,-2,-1,1,2,3,4} \grady{1,2,3}
  \draw[coulCourbe,line width=1.1pt,domain=-4:4,samples=60] plot (\x,{exp(0.3*\x)});
  \draw[coulCourbeB,line width=1.1pt,domain=-4:4,samples=60] plot (\x,{exp(-0.3*\x)});
  \node[coulCourbe,font=\footnotesize] at (4.4,3.1) {$\Cf$};
  \node[coulCourbeB,font=\footnotesize] at (-4.4,3.1) {$\Cg$};
\end{tikzpicture}
\end{center}
\end{minipage}
\cas{4.}
\explique{\res{$f$ modélise une croissance exponentielle, et $g$ une décroissance
exponentielle.}}{$\e^{kt}$ et $\e^{-kt}$ avec $k>0$}
\end{exo}

\begin{exo}[Reconnaître une suite géométrique]
\cas{a)}
\begin{calculs}
u_{n+1} &= \e^{3(n+1)} &&\justif{on remplace $n$ par $n+1$}\\
u_{n+1} &= \e^{3n}\times\e^{3} &&\justif{$\e^{a+b}=\e^{a}\times\e^{b}$}\\
u_{n+1} &= u_{n}\times\e^{3} &&
\end{calculs}
\explique{$(u_{n})$ est géométrique de \res{raison $\e^{3}$ et de premier terme $u_{0}=1$} ;
$\e^{3}>1$ et $u_{0}>0$ : \res{elle est croissante}.}{$u_{0}=\e^{0}$}
\cas{b)}
\begin{calculs}
v_{n+1} &= 4\,\e^{-0{,}5(n+1)} &&\\
v_{n+1} &= 4\,\e^{-0{,}5n}\times\e^{-0{,}5} &&\\
v_{n+1} &= v_{n}\times\e^{-0{,}5} &&
\end{calculs}
\explique{$(v_{n})$ est géométrique de \res{raison $\e^{-0{,}5}$ et de premier terme $v_{0}=4$} ;
$0<\e^{-0{,}5}<1$ et $v_{0}>0$ : \res{elle est décroissante}.}{$\e^{-0{,}5}\approx 0{,}61$}
\cas{c)}
\begin{calculs}
w_{n+1} &= \e^{2(n+1)+1} &&\\
w_{n+1} &= \e^{2n+1}\times\e^{2} &&\justif{$2n+3=(2n+1)+2$}\\
w_{n+1} &= w_{n}\times\e^{2} &&
\end{calculs}
\explique{$(w_{n})$ est géométrique de \res{raison $\e^{2}$ et de premier terme $w_{0}=\e$} ;
$\e^{2}>1$ et $w_{0}>0$ : \res{elle est croissante}.}{$w_{0}=\e^{1}$}
\end{exo}

\begin{exo}[Élimination d'un médicament]
\cas{1.}
\begin{calculs}
Q(0) &= 80\,\e^{0}=\res{80} &&\justif{$\e^{0}=1$}
\end{calculs}
\explique{\res{Juste après la prise, l'organisme contient $80$~mg de médicament.}}{on
interprète}
\cas{2.}
\begin{calculs}
Q'(t) &= 80\times(-0{,}1)\,\e^{-0{,}1t} &&\justif{$\left(\e^{at}\right)'=a\,\e^{at}$}\\
Q'(t) &= -8\,\e^{-0{,}1t} &&
\end{calculs}
\explique{$Q'(t)<0$ : $Q$ est \res{strictement décroissante}.}{$\e^{-0{,}1t}>0$}
\cas{3.}
\begin{calculs}
u_{n+1} &= 80\,\e^{-0{,}1(n+1)} &&\\
u_{n+1} &= 80\,\e^{-0{,}1n}\times\e^{-0{,}1} &&\justif{$\e^{a+b}=\e^{a}\times\e^{b}$}\\
u_{n+1} &= u_{n}\times\e^{-0{,}1} &&
\end{calculs}
\explique{$(u_{n})$ est géométrique de raison \res{$\e^{-0{,}1}\approx 0{,}905$}.}{premier terme
$u_{0}=80$}
\cas{4.}
\explique{$0{,}905=1-0{,}095$ : la quantité baisse d'environ \res{$9{,}5\,\%$} par heure.}{coefficient
multiplicateur $1+t$, avec $t=-0{,}095$}
\cas{5.}
\explique{À la calculatrice, $Q(13)\approx 21{,}8$ et $Q(14)\approx 19{,}7$ ; $Q$ est
décroissante.}{la quantité reste ensuite sous $20$~mg}
\explique{Il reste moins de $20$~mg \res{au bout de $14$~heures}.}{}
\end{exo}

\begin{exo}[Du modèle discret au modèle continu]
\cas{1.}
\explique{Une hausse de $4\,\%$ multiplie la population par $1{,}04$ chaque année :
$u_{n+1}=1{,}04\,u_{n}$.}{coefficient multiplicateur $1+0{,}04$}
\explique{$(u_{n})$ est géométrique de raison $1{,}04$ et de premier terme $u_{0}=2\,000$, donc
\res{$u_{n}=2\,000\times 1{,}04^{n}$}.}{$u_{n}=u_{0}\,q^{n}$}
\cas{2.}
\begin{calculs}
f(1) &= 2\,000\,\e^{k} &&\\
u_{1} &= 2\,000\times 1{,}04 &&\justif{on veut $f(1)=u_{1}$}
\end{calculs}
\explique{On divise par $2\,000$ : \res{$\e^{k}=1{,}04$}.}{la raison de la suite est $\e^{k}$
(propriété 8)}
\cas{3.}
\explique{À la calculatrice, $\e^{0{,}039}\approx 1{,}0398$ et $\e^{0{,}040}\approx 1{,}0408$ :
$1{,}04$ est plus proche de la première valeur, donc \res{$k\approx 0{,}039$}.}{tableau de
valeurs de $\e^{x}$ au pas de $0{,}001$}
\cas{4.}
\begin{calculs}
f(2{,}5) &= 2\,000\,\e^{0{,}039\times 2{,}5} &&\\
f(2{,}5) &= 2\,000\,\e^{0{,}0975} &&\\
f(2{,}5) &\approx \res{2\,205} &&\justif{à la calculatrice}
\end{calculs}
\explique{$u_{2}=2\,163{,}2$ et $u_{3}\approx 2\,249{,}7$ : \res{la population au milieu de 2028
est comprise entre celles de 2028 et de 2029}, environ $2\,205$ habitants.}{le modèle
continu donne des valeurs entre deux années}
\end{exo}

\partie{Problèmes concrets}

\begin{exo}[Bénéfice d'un atelier]
\cas{1.}
\begin{calculs}
B(1) &= (4-20)\,\e^{-0{,}2} &&\\
B(1) &= -16\,\e^{-0{,}2} &&\\
B(1) &\approx \res{-13{,}1} &&\justif{à la calculatrice}
\end{calculs}
\explique{\res{Pour $100$ trottinettes, l'atelier perd environ $13\,100$~€.}}{$B(x)$ est en
milliers d'euros}
\cas{2.}
\explique{$\e^{-0{,}2x}>0$ : $B(x)\ge 0$ quand $4x-20\ge 0$, soit $x\ge 5$ ;
$S=\res{\intff{5}{20}}$.}{produit par une exponentielle (exercice 9)}
\explique{\res{L'atelier ne perd pas d'argent à partir de $500$ trottinettes.}}{$x$ en
centaines}
\cas{3.}
\begin{calculs}
u(x) &= 4x-20 &&\justif{$u'(x)=4$}\\
v(x) &= \e^{-0{,}2x} &&\justif{$v'(x)=-0{,}2\,\e^{-0{,}2x}$}\\
B'(x) &= 4\,\e^{-0{,}2x}+(4x-20)\times\left(-0{,}2\,\e^{-0{,}2x}\right) &&\justif{$u'v+uv'$}\\
B'(x) &= (4-0{,}8x+4)\,\e^{-0{,}2x} &&\justif{$-0{,}2\times(4x-20)=-0{,}8x+4$}\\
B'(x) &= (8-0{,}8x)\,\e^{-0{,}2x} &&
\end{calculs}
\cas{4.}
\explique{$\e^{-0{,}2x}>0$ : $B'(x)$ a le signe de $8-0{,}8x$, qui s'annule en $x=10$, positif
avant.}{$8-0{,}8x>0$ pour $x<10$}
\begin{calculs}
B(10) &= 20\,\e^{-2} &&\justif{environ $2{,}707$}\\
B(20) &= 60\,\e^{-4} &&\justif{environ $1{,}099$}
\end{calculs}
\begin{center}
\begin{tikzpicture}[scale=0.9]
\tkzTabInit[lgt=2,espcl=2.6]{$x$/0.8,$B'(x)$/0.8,$B(x)$/1.6}{$1$,$10$,$20$}
\tkzTabLine{,+,z,-,}
\tkzTabVar{-/$-16\e^{-0{,}2}$,+/$20\e^{-2}$,-/$60\e^{-4}$}
\end{tikzpicture}
\end{center}
\explique{$B$ est \res{croissante sur $\intff{1}{10}$, décroissante sur $\intff{10}{20}$}.}{}
\cas{5.}
\explique{Le maximum est atteint pour $x=10$ : il faut fabriquer \res{$1\,000$ trottinettes},
pour un bénéfice de $20\,\e^{-2}\approx 2{,}70671$ milliers d'euros, soit
\res{$2\,707$~€}.}{arrondi à l'euro}
\end{exo}

\begin{exo}[Refroidissement d'un four à pain]
\cas{1.}
\begin{calculs}
T(0) &= 20+780\,\e^{0} &&\\
T(0) &= \res{800} &&
\end{calculs}
\explique{\res{Quand on l'éteint, le four est à $800$~\textdegree C.}}{}
\cas{2.}
\begin{calculs}
T'(t) &= 780\times(-0{,}1)\,\e^{-0{,}1t} &&\justif{la constante $20$ a pour dérivée $0$}\\
T'(t) &= \res{-78\,\e^{-0{,}1t}} &&
\end{calculs}
\explique{$T'(t)<0$ : $T$ est \res{strictement décroissante} ; le four refroidit.}{$\e^{-0{,}1t}>0$}
\cas{3.}
\begin{calculs}
T'(0) &= -78\,\e^{0}=\res{-78} &&
\end{calculs}
\explique{\res{Juste après l'extinction, la température baisse de $78$~\textdegree C par heure.}}{$T'(t)$
est une vitesse, en \textdegree C par heure}
\cas{4.}
\explique{$T(22)\approx 106{,}4$ est au-dessus de $100$, $T(23)\approx 98{,}2$ en dessous, et $T$
est décroissante : la température reste ensuite sous $100$~\textdegree C.}{on lit le tableau}
\explique{On peut nettoyer le four \res{$23$~heures après l'extinction}, soit à \res{21~h le
lendemain}.}{$22$~h $+\,23$~h}
\end{exo}

\begin{exo}[Le plus grand potager]
\cas{1.}
\explique{$OA=x$ et $AB=f(x)=6\,\e^{-0{,}5x}$, car $B$ est le point de $\Cf$
d'abscisse $x$.}{longueur et largeur du rectangle}
\begin{calculs}
g(x) &= OA\times AB &&\\
g(x) &= x\times 6\,\e^{-0{,}5x} &&\\
g(x) &= \res{6x\,\e^{-0{,}5x}} &&
\end{calculs}
\cas{2.}
\begin{calculs}
u(x) &= 6x &&\justif{$u'(x)=6$}\\
v(x) &= \e^{-0{,}5x} &&\justif{$v'(x)=-0{,}5\,\e^{-0{,}5x}$}\\
g'(x) &= 6\,\e^{-0{,}5x}+6x\times\left(-0{,}5\,\e^{-0{,}5x}\right) &&\justif{$u'v+uv'$}\\
g'(x) &= (6-3x)\,\e^{-0{,}5x} &&\justif{on factorise par $\e^{-0{,}5x}$}
\end{calculs}
\cas{3.}
\explique{$\e^{-0{,}5x}>0$ : $g'(x)$ a le signe de $6-3x$, positif pour $x<2$, négatif
après.}{$6-3x=0$ pour $x=2$}
\begin{calculs}
g(0) &= 0 &&\\
g(2) &= 12\,\e^{-1} &&\justif{environ $4{,}41$}\\
g(5) &= 30\,\e^{-2{,}5} &&\justif{environ $2{,}46$}
\end{calculs}
\begin{center}
\begin{tikzpicture}[scale=0.9]
\tkzTabInit[lgt=2,espcl=2.6]{$x$/0.8,$g'(x)$/0.8,$g(x)$/1.6}{$0$,$2$,$5$}
\tkzTabLine{,+,z,-,}
\tkzTabVar{-/$0$,+/$12\e^{-1}$,-/$30\e^{-2{,}5}$}
\end{tikzpicture}
\end{center}
\explique{$g$ est \res{croissante sur $\intff{0}{2}$, décroissante sur $\intff{2}{5}$}.}{}
\cas{4.}
\explique{Le maximum de $g$ est atteint pour $x=2$ : on place \res{$A$ à $2$~m de $O$}.
L'aire est alors $12\,\e^{-1}\approx 4{,}4146$~m$^{2}$, soit \res{$4{,}41$~m$^{2}$} au dm$^{2}$
près.}{$1$~dm$^{2}=0{,}01$~m$^{2}$}
\end{exo}

\end{document}
